\documentclass[../../main.tex]{subfiles}

\begin{document}

\section{测度与可测集}\label{section:ch17-measurable-sets}

回忆一个集合 $X$ 的子集的 $\sigma$ 代数是 $X$ 的一个子集族, 它包含空集且对 $X$ 的补以及可数并封闭, 根据 De Morgan 等式, 它关于交也是封闭的. 谈到集函数 $\mu$ 我们指的是赋予某些集一个扩充的实数的函数.

\begin{definition}[可测空间与可测集]\label{definition:ch17-measurable-space}
谈到可测空间 $(X,\mathcal{M})$, 我们指的是由集合 $X$ 与 $X$ 的子集的 $\sigma$ 代数 $\mathcal{M}$ 组成的一对 $(X,\mathcal{M})$. 若 $X$ 的子集 $E$ 属于 $\mathcal{M}$, 则 $E$ 称为可测的(或关于 $\mathcal{M}$ 可测).
\end{definition}

\begin{definition}[测度与测度空间]\label{definition:ch17-measure}
谈到可测空间 $(X,\mathcal{M})$ 上的测度, 我们指的是扩充实值非负函数 $\mu:\mathcal{M}\to[0,\infty]$, 使得 $\mu(\varnothing)=0$ 且在以下意义下是可数可加的: 对任何可测集的可数不交族 $\{E_{k}\}_{k=1}^{\infty}$,
\begin{align*}
\mu\left(\bigcup_{k=1}^{\infty}E_{k}\right)=\sum_{k=1}^{\infty}\mu(E_{k}).
\end{align*}
谈到测度空间 $(X,\mathcal{M},\mu)$, 我们指的是可测空间 $(X,\mathcal{M})$ 以及定义在 $\mathcal{M}$ 上的测度 $\mu$.
\end{definition}

测度空间的一个例子是 $(\mathbf{R},\mathcal{L},m)$, 其中 $\mathbf{R}$ 是实数集, $\mathcal{L}$ 是实数的 Lebesgue 可测集族, $m$ 是 Lebesgue 测度. 测度空间的第二个例子是 $(\mathbf{R},\mathcal{B},m)$, 其中 $\mathcal{B}$ 是实数的 Borel 集族, 而 $m$ 仍然是 Lebesgue 测度. 对任何集合 $X$, 我们定义 $\mathcal{M}=2^{X}$, 即 $X$ 的所有子集组成的族. 通过定义有限集的测度为该集中元素的个数而无穷集的测度为 $\infty$, 定义了测度 $\eta$. 我们称 $\eta$ 为 $X$ 上的计数测度. 对任何 $X$ 的子集的任何 $\sigma$ 代数 $\mathcal{M}$ 和属于 $X$ 的点 $x_{0}$, 集中在 $x_{0}$ 的 Dirac 测度, 记为 $\delta_{x_{0}}$, 赋予 $\mathcal{M}$ 中包含 $x_{0}$ 的集测度为 $1$ 而不包含 $x_{0}$ 的集测度为 $0$: 这定义了 Dirac 测度空间 $(X,\mathcal{M},\delta_{x_{0}})$. 以下是有点奇特的例子: 令 $X$ 为任何不可数集, 而 $\mathcal{C}$ 是 $X$ 的那些或者是可数集或者是可数集的补集的子集构成的族. 则 $\mathcal{C}$ 是 $\sigma$ 代数且通过对 $X$ 的每个可数子集 $A$ 设 $\mu(A)=0$ 而对 $X$ 的每个补集是可数的子集 $B$ 设 $\mu(B)=1$, 我们定义了 $\mathcal{C}$ 上的测度. 则 $(X,\mathcal{C},\mu)$ 是一个测度空间.

观察到以下性质是有用的. 对任何测度空间 $(X,\mathcal{M},\mu)$, 若 $X_{0}$ 属于 $\mathcal{M}$, 则 $(X_{0},\mathcal{M}_{0},\mu_{0})$ 也是一个测度空间, 其中 $\mathcal{M}_{0}$ 是 $\mathcal{M}$ 的包含于 $X_{0}$ 的子集族而 $\mu_{0}$ 是 $\mu$ 在 $\mathcal{M}_{0}$ 上的限制.

\begin{proposition}\label{proposition:ch17-basic-properties}
令 $(X,\mathcal{M},\mu)$ 为测度空间.
\begin{description}
\item[（有限可加性）] 对于可测集的任何有限不交族 $\{E_{k}\}_{k=1}^{n}$,
\begin{align*}
\mu\left(\bigcup_{k=1}^{n}E_{k}\right)=\sum_{k=1}^{n}\mu(E_{k}).
\end{align*}
\item[（单调性）] 若 $A$ 和 $B$ 是可测集且 $A\subseteq B$, 则
\begin{align*}
\mu(A)\leqslant\mu(B).
\end{align*}
\item[（分割性）] 此外, 若 $A\subseteq B$ 且 $\mu(A)<\infty$, 则
\begin{align*}
\mu(B\sim A)=\mu(B)-\mu(A).
\end{align*}
因此, 若 $\mu(A)=0$, 则
\begin{align*}
\mu(B\sim A)=\mu(B).
\end{align*}
\item[（可数单调性）] 对任何覆盖可测集 $E$ 的可测集族 $\{E_{k}\}_{k=1}^{\infty}$,
\begin{align*}
\mu(E)\leqslant\sum_{k=1}^{\infty}\mu(E_{k}).
\end{align*}
\end{description}
\end{proposition}

\begin{proof}
通过对 $k>n$ 设 $E_{k}=\varnothing$, 使得 $\mu(E_{k})=0$, 有限可加性从可数可加性得到. 根据有限可加性,
\begin{align*}
\mu(B)=\mu(A)+\mu(B\sim A)
\end{align*}
立即推出单调性与分割性. 为证明可数单调性, 定义 $G_{1}=E_{1}$ 且对所有 $k\geqslant 2$, 定义 $G_{k}=E_{k}\sim\left[\bigcup_{i=1}^{k-1}E_{i}\right]$. 观察到
\begin{align*}
\{G_{k}\}_{k=1}^{\infty}\text{ 是不交的, }\bigcup_{k=1}^{\infty}G_{k}=\bigcup_{k=1}^{\infty}E_{k},\text{ 且对所有 }k,G_{k}\subseteq E_{k}.
\end{align*}
从 $\mu$ 的单调性与可数可加性我们推出
\begin{align*}
\mu(E)\leqslant\mu\left(\bigcup_{k=1}^{\infty}E_{k}\right)=\mu\left(\bigcup_{k=1}^{\infty}G_{k}\right)=\sum_{k=1}^{\infty}\mu(G_{k})\leqslant\sum_{k=1}^{\infty}\mu(E_{k}).
\end{align*}
\end{proof}

可数单调性是可数可加性与单调性的混合, 它频繁地被援引, 故我们单独对它命名.

集合序列 $\{E_{k}\}_{k=1}^{\infty}$ 称为上升的, 若对每个 $k$, $E_{k}\subseteq E_{k+1}$; 称为下降的, 若对每个 $k$, $E_{k+1}\subseteq E_{k}$.

\begin{proposition}[测度的连续性]\label{proposition:ch17-continuity}
令 $(X,\mathcal{M},\mu)$ 为测度空间.
\begin{description}
\item[（i）] 若 $\{A_{k}\}_{k=1}^{\infty}$ 是上升的可测集序列, 则
\begin{align}
\mu\left(\bigcup_{k=1}^{\infty}A_{k}\right)=\lim_{k\to\infty}\mu(A_{k}).\label{align:ch17-cont-up}
\end{align}
\item[（ii）] 若 $\{B_{k}\}_{k=1}^{\infty}$ 是满足 $\mu(B_{1})<\infty$ 的下降的可测集序列, 则
\begin{align}
\mu\left(\bigcap_{k=1}^{\infty}B_{k}\right)=\lim_{k\to\infty}\mu(B_{k}).\label{align:ch17-cont-down}
\end{align}
\end{description}
\end{proposition}

测度的连续性的证明与实直线上的 Lebesgue 测度的连续性的证明逐字相同, 见 2.5 节定理 15.

对于一个测度空间 $(X,\mathcal{M},\mu)$ 和 $X$ 的一个可测子集 $E$, 我们说一个性质在 $E$ 上几乎处处成立或对 $E$ 中的几乎所有 $x$ 成立, 若它在 $E\sim E_{0}$ 上成立, 其中 $E_{0}$ 是 $E$ 的可测子集且 $\mu(E_{0})=0$.

\begin{lemma}[Borel-Cantelli]\label{lemma:ch17-borel-cantelli}
令 $(X,\mathcal{M},\mu)$ 为测度空间而 $\{E_{k}\}_{k=1}^{\infty}$ 是满足 $\sum_{k=1}^{\infty}\mu(E_{k})<\infty$ 的可数可测集族. 则 $X$ 中的几乎所有 $x$ 属于至多有限个 $E_{k}$.
\end{lemma}

\begin{proof}
对每个 $n$, 根据 $\mu$ 的可数单调性, $\mu\left(\bigcup_{k=n}^{\infty}E_{k}\right)\leqslant\sum_{k=n}^{\infty}\mu(E_{k})$. 因此根据 $\mu$ 的连续性,
\begin{align*}
\mu\left(\bigcap_{n=1}^{\infty}\left[\bigcup_{k=n}^{\infty}E_{k}\right]\right)=\lim_{n\to\infty}\mu\left(\bigcup_{k=n}^{\infty}E_{k}\right)\leqslant\lim_{n\to\infty}\sum_{k=n}^{\infty}\mu(E_{k})=0.
\end{align*}
观察到 $\bigcap_{n=1}^{\infty}\left[\bigcup_{k=n}^{\infty}E_{k}\right]$ 是 $X$ 中属于无穷多个 $E_{k}$ 的点集.
\end{proof}

\begin{definition}[有限与 $\sigma$ 有限测度]\label{definition:ch17-finite-sigma-finite}
令 $(X,\mathcal{M},\mu)$ 为测度空间. 测度 $\mu$ 称为有限的, 若 $\mu(X)<\infty$. 它称为 $\sigma$ 有限的, 若 $X$ 是可数个有限测度集的并. 可测集 $E$ 称为具有有限测度, 若 $\mu(E)<\infty$; 称为是 $\sigma$ 有限的, 若 $E$ 是可测集的可数族的并, 其中该族的每个集合有有限测度.
\end{definition}

考虑到 $\sigma$ 有限的准则, 由有限测度集组成的可数覆盖可取为不交的. 事实上, 若 $\{X_{k}\}_{k=1}^{\infty}$ 是这样一个覆盖, 对 $k\geqslant 2$, 用 $X_{k}\sim\bigcup_{j=1}^{k-1}X_{j}$ 代替每个 $X_{k}$ 得到了由有限测度集组成的不交覆盖. $[0,1]$ 上的 Lebesgue 测度是有限测度的一个例子, 而 $(-\infty,\infty)$ 上的 Lebesgue 测度是 $\sigma$ 有限测度的一个例子. 不可数集上的计数测度不是 $\sigma$ 有限的.

实直线上的 Lebesgue 测度的熟悉性质以及一元实变量的函数的 Lebesgue 积分对任意 $\sigma$ 有限测度成立. 抽象测度论的许多处理中限定为 $\sigma$ 有限测度. 然而, 一般理论的许多部分不要求 $\sigma$ 有限的假设, 看起来在这个不必要的限制上建立理论不令人满意.

\begin{definition}[完备测度空间]\label{definition:ch17-complete}
测度空间 $(X,\mathcal{M},\mu)$ 称为完备的, 若 $\mathcal{M}$ 包含所有测度为零的子集, 即若 $E$ 属于 $\mathcal{M}$ 且 $\mu(E)=0$, 则 $E$ 的每个子集也属于 $\mathcal{M}$.
\end{definition}

我们证明了实直线上的 Lebesgue 测度是完备的. 此外, 我们也证明了 Cantor 集, 作为一个 Lebesgue 测度为零的 Borel 集, 包含不是 Borel 集的子集, 见 2.7 节命题 22. 因此限定在 Borel 集的 $\sigma$ 代数的实直线上的 Lebesgue 测度不是完备的. 以下命题的证明留给读者(习题 9), 它告诉我们每个测度空间可以完备化. 该命题中描述的测度空间 $(X,\mathcal{M}_{0},\mu_{0})$ 称为 $(X,\mathcal{M},\mu)$ 的完备化.

\begin{proposition}\label{proposition:ch17-completion}
令 $(X,\mathcal{M},\mu)$ 为测度空间. 定义 $\mathcal{M}_{0}$ 为 $X$ 的形如 $E=A\cup B$ 的子集族, 其中 $B\in\mathcal{M}$, $A\subseteq C$, 这里 $C\in\mathcal{M}$, 且满足 $\mu(C)=0$. 对这样的 $E$ 定义 $\mu_{0}(E)=\mu(B)$. 则 $\mathcal{M}_{0}$ 是包含 $\mathcal{M}$ 的 $\sigma$ 代数, $\mu_{0}$ 是 $\mu$ 的延拓, 而 $(X,\mathcal{M}_{0},\mu_{0})$ 是完备测度空间.
\end{proposition}

\begin{exercise}\label{exercise:ch17-01}
令 $f$ 为 $\mathbf{R}$ 上的非负 Lebesgue 可测函数, 对 $\mathbf{R}$ 的每个 Lebesgue 可测子集 $E$, 定义 $\mu(E)=\int_{E}f$ 为 $f$ 在 $E$ 上的 Lebesgue 积分. 证明 $\mu$ 是 $\mathbf{R}$ 的 Lebesgue 可测子集的 $\sigma$ 代数上的测度.
\end{exercise}

\begin{exercise}\label{exercise:ch17-02}
令 $\mathcal{M}$ 是 $X$ 的子集的 $\sigma$ 代数, 而集函数 $\mu:\mathcal{M}\to[0,\infty)$ 是有限可加的. 证明 $\mu$ 是测度当且仅当只要 $\{A_{k}\}_{k=1}^{\infty}$ 是上升的序列就有
\begin{align*}
\mu\left(\bigcup_{k=1}^{\infty}A_{k}\right)=\lim_{k\to\infty}\mu(A_{k}).
\end{align*}
\end{exercise}

\begin{exercise}\label{exercise:ch17-03}
令 $\mathcal{M}$ 为集合 $X$ 的子集的 $\sigma$ 代数. 对 $\mathcal{M}$ 中集合的递减序列, 叙述与证明前一个习题的相应结论.
\end{exercise}

\begin{exercise}\label{exercise:ch17-04}
令 $\{(X_{\lambda},\mathcal{M}_{\lambda},\mu_{\lambda})\}_{\lambda\in\Lambda}$ 为由集合 $\Lambda$ 参数化的测度空间族. 假设集族 $\{X_{\lambda}\}_{\lambda\in\Lambda}$ 是不交的. 我们可以构建新的测度空间 $(X,\mathcal{B},\mu)$(称为 $\{(X_{\lambda},\mathcal{M}_{\lambda},\mu_{\lambda})\}_{\lambda\in\Lambda}$ 的并), 方法是: 令 $X=\bigcup_{\lambda\in\Lambda}X_{\lambda}$, $\mathcal{B}$ 为 $X$ 的子集 $B$ 构成的族, 满足对每个 $\lambda\in\Lambda$, $B\cap X_{\lambda}\in\mathcal{M}_{\lambda}$, 且对 $B\in\mathcal{B}$ 定义 $\mu(B)=\sum_{\lambda\in\Lambda}\mu_{\lambda}[B\cap X_{\lambda}]$.
\begin{description}
\item[（i）] 证明 $\mathcal{B}$ 是 $\sigma$ 代数.
\item[（ii）] 证明 $\mu$ 是 $\sigma$ 测度.
\item[（iii）] 证明 $\mu$ 是 $\sigma$ 有限的当且仅当除可数个测度 $\mu_{\lambda}$ 外, 有 $\mu(X_{\lambda})=0$ 且其余测度是 $\sigma$ 有限的.
\end{description}
\end{exercise}

\begin{exercise}\label{exercise:ch17-05}
令 $(X,\mathcal{M},\mu)$ 为测度空间. $X$ 的两个子集 $E_{1}$ 和 $E_{2}$ 的对称差 $E_{1}\Delta E_{2}$ 定义为
\begin{align*}
E_{1}\Delta E_{2}=[E_{1}\sim E_{2}]\cup[E_{2}\sim E_{1}].
\end{align*}
\begin{description}
\item[（i）] 证明: 若 $E_{1}$ 和 $E_{2}$ 是可测的且 $\mu(E_{1}\Delta E_{2})=0$, 则 $\mu(E_{1})=\mu(E_{2})$.
\item[（ii）] 证明: 若 $\mu$ 是完备的且 $E_{1}\in\mathcal{M}$, 则 $E_{2}\in\mathcal{M}$, 若 $\mu(E_{1}\Delta E_{2})=0$.
\end{description}
\end{exercise}

\begin{exercise}\label{exercise:ch17-06}
令 $(X,\mathcal{M},\mu)$ 为测度空间, $X_{0}$ 属于 $\mathcal{M}$. 定义 $\mathcal{M}_{0}$ 为 $\mathcal{M}$ 中 $X_{0}$ 的子集族, 而 $\mu_{0}$ 为 $\mu$ 在 $\mathcal{M}_{0}$ 上的限制. 证明 $(X_{0},\mathcal{M}_{0},\mu_{0})$ 是测度空间.
\end{exercise}

\begin{exercise}\label{exercise:ch17-07}
令 $(X,\mathcal{M})$ 为可测空间. 证明:
\begin{description}
\item[（i）] 若 $\mu$ 和 $\nu$ 是定义在 $\mathcal{M}$ 上的测度, 则 $\mathcal{M}$ 上定义为 $\lambda(E)=\mu(E)+\nu(E)$ 的集函数 $\lambda$ 也是一个测度. 我们记 $\lambda$ 为 $\mu+\nu$.
\item[（ii）] 若 $\mu$ 和 $\nu$ 是定义在 $\mathcal{M}$ 上的测度且 $\mu\geqslant\nu$, 则存在 $\mathcal{M}$ 上的测度 $\lambda$ 使得 $\mu=\nu+\lambda$.
\item[（iii）] 若 $\nu$ 是 $\sigma$ 有限的, 则（ii）中的测度 $\lambda$ 是唯一的.
\item[（iv）] $\lambda$ 一般不唯一但总存在最小的 $\lambda$.
\end{description}
\end{exercise}

\begin{exercise}\label{exercise:ch17-08}
令 $(X,\mathcal{M},\mu)$ 为测度空间. 测度 $\mu$ 称为是半有限的, 若每个具有无穷测度的可测集包含任意大的有限测度的可测集.
\begin{description}
\item[（i）] 证明每个 $\sigma$ 有限测度是半有限的.
\item[（ii）] 对 $E\in\mathcal{M}$, 定义 $\mu_{1}(E)=\sup\{\mu(F)\mid F\subseteq E,\mu(F)<\infty\}$. 证明 $\mu_{1}$ 是半有限测度: 它称为 $\mu$ 的半有限部分.
\item[（iii）] 找到 $\mathcal{M}$ 上的测度 $\mu_{2}$, 它仅取 $0$ 与 $\infty$ 值且 $\mu=\mu_{1}+\mu_{2}$.
\end{description}
\end{exercise}

\begin{exercise}\label{exercise:ch17-09}
证明命题 3, 即证明 $\mathcal{M}_{0}$ 是 $\sigma$ 代数, $\mu_{0}$ 是恰当定义的, 而 $(X,\mathcal{M}_{0},\mu_{0})$ 是完备的. 在什么意义下 $\mathcal{M}_{0}$ 是最小的?
\end{exercise}

\begin{exercise}\label{exercise:ch17-10}
若 $(X,\mathcal{M},\mu)$ 是测度空间, 我们说 $X$ 的子集 $E$ 是局部可测的, 若对每个满足 $\mu(B)<\infty$ 的 $B\in\mathcal{M}$, 交 $E\cap B$ 属于 $\mathcal{M}$. 测度 $\mu$ 称为饱和的, 若每个局部可测集是可测的.
\begin{description}
\item[（i）] 证明每个 $\sigma$ 有限测度是饱和的.
\item[（ii）] 证明局部可测集族 $\mathcal{C}$ 是 $\sigma$ 代数.
\item[（iii）] 令 $(X,\mathcal{M},\mu)$ 为测度空间而 $\mathcal{C}$ 是局部可测集的 $\sigma$ 代数. 对 $E\in\mathcal{C}$, 定义 $\overline{\mu}(E)=\mu(E)$, 若 $E\in\mathcal{M}$; 而 $\overline{\mu}(E)=\infty$, 若 $E\notin\mathcal{M}$. 证明 $(X,\mathcal{C},\overline{\mu})$ 是饱和的测度空间.
\item[（iv）] 若 $\mu$ 是半有限的且 $E\in\mathcal{C}$, 定义 $\underline{\mu}(E)=\sup\{\mu(B)\mid B\in\mathcal{M},B\subseteq E\}$. 证明 $(X,\mathcal{C},\underline{\mu})$ 是饱和的测度空间而 $\underline{\mu}$ 是 $\mu$ 的延拓. 给出一个例子说明 $\overline{\mu}$ 和 $\underline{\mu}$ 可以不同.
\end{description}
\end{exercise}

\begin{exercise}\label{exercise:ch17-11}
令 $\mu$ 和 $\eta$ 为可测空间 $(X,\mathcal{M})$ 上的测度. 对 $E\in\mathcal{M}$, 定义 $\nu(E)=\max\{\mu(E),\eta(E)\}$. $\nu$ 是 $(X,\mathcal{M})$ 上的测度吗?
\end{exercise}

\end{document}